What is the actuarial equivalence principle / principle of pooling? 
What are systematic risks?


The most important principle of the insurance business is the pooling of risks:
- Many homogeneous and independent risks in the insurer’s portfolio
  fair premium=average claims amount (plus safety loading, expenses)
Premium calculation based on the Actuarial Equivalence Principle:
- Expected present value of premium payments=expected present value of insurance benefits
- Pricing of insurance contracts related to statistical probabilities (‘statistical’ or ‘real-world’ probability measure \( \mathbb{P}\))
Insurance pricing considers non-systematic risks:
For systematic risk like financial market risk and non-diversifiable insurance risks (e. g., longevity risk) the
equivalence principle is not applicable, due to the lack of pooling of risks.

Why is the net risk premium a good approximation, but not fully sufficient? One
static perspective is the CLT. Another, dynamic one uses the law of iterated logarithm.

Pooling of risks expressed in terms of the Law of Large Numbers:
- For a large portfolio of i.i.d. insurance claims \(L_i \sim L\), \(i \in \mathbb{N}\), 
   modeled as non-negative random variables on \((\Omega, \mathcal{F}, P)\),
  \[\lim_{n \rightarrow \infty} \frac{1}{n} \sum_{i=1}^n L_i = E[L] \quad P\text{-a.s..}\]
- This suggests that the premium per contract should be equal to the expected claims amount ('net risk premium'), i.e., \(\pi(L) = E[L]\).

ATTENTION: The net risk premium is not sufficient:
- Classical results of ruin theory state that ruin of the insurer occurs in the long run \(P\)-a.s. if only the net risk premium \(\pi(L) = E[L]\) is charged.
- A similar result holds in a one-period setting: Letting the portfolio size \(n\) tend to infinity, the one-period 'ruin probability'
\[P \left[ \pi n(L) - \sum_{i=1}^n L_i < 0 \right] = P \left[ \frac{\sum_{i=1}^n (L_i - E[L])}{\sqrt{n} \sigma(L)} < 0 \right]\]
converges to 1/2 (Central Limit Theorem).

Conclusion: An (explicit or implicit) safety loading is necessary!
By the strong law of large numbers, the random walk \(Z_n := \sum_{i=1}^n (\pi + L_i)\), \(n \in \mathbb{N}\), satisfies
\[\lim_{n \rightarrow \infty} \frac{1}{n} Z_n = \pi + E[L_1] \quad P\text{-a.s..}\]
\(\pi < E[L_1]\): For any \(u\), this implies
\[\lim_{n \rightarrow \infty} K_n^u = \lim_{n \rightarrow \infty} Z_n = -\infty \quad P\text{-a.s.},\]
hence \(P[\inf_{n \in \mathbb{N}} K_n^u < 0] = 1\).
\(\pi > E[L_1]\): In this case, \(\lim_{n \rightarrow \infty} Z_n = +\infty \quad P\text{-a.s.}\) and hence
\[\inf_{n \in \mathbb{N}} Z_n > -\infty \quad P\text{-a.s..}\]
This implies
\[\lim_{u \rightarrow \infty} P[\inf_{n \in \mathbb{N}} K_n^u < 0] = \lim_{u \rightarrow \infty} P[u + \inf_{n \in \mathbb{N}} Z_n < 0] = 0.\]

\(\pi = E[L_1]\):
For \(\pi = E[L_1]\), the increments of the random walk \(Z_n\), \(n \in \mathbb{N}\), have zero expectation.
The law of the iterated logarithm yields
\[\lim_{n \rightarrow \infty} \frac{Z_n}{\sqrt{2nVar[L_1] \ln \ln n}} = -1 \quad P\text{-a.s.},\]
i.e., for \(P\)-almost all \(\omega \in \Omega\) there exists a subsequence \(n_k(\omega)\) such that
\[\lim_{k \rightarrow \infty} Z_{n_k(\omega)}(\omega) = -\infty.\]
This implies as above
\[P[\inf_{n \in \mathbb{N}} K_n^u < 0] = P[u + \inf_{n \in \mathbb{N}} Z_n < 0] = 1.\]


Explain the terms explicit premium principle and implicit premium principle. How
can the latter be framed using the notion of first order stochastic dominance?

Example 4.7.1 (Explicit premium principle)

Let \(L\) be an insurance claim, and let \(\rho(-L)\) be some measure of risk (e.g., variance, standard deviation, Value at Risk, ...). 
For a parameter \(a>0\), the functional
\[\pi(L) = E[L] + a\rho(-L)\]
is called an explicit premium principle.

Remarks:
Safety loadings might also be chosen dependent on the number of risks in a portfolio to keep the probability of ruin below a given threshold.

An implicit safety loading results from conservative best estimate assumptions:
Assume that \(F\) is the cumulative distribution function of a random variable \(L \geq 0\) under \(\mathbb{P}\), 
but the actuary calculates with a more conservative cumulative distribution function \(G\) (corresponding to a measure \(\tilde{\mathbb{P}}\)), 
i.e.,
\[1 - F(x) = \mathbb{P}[L>x] \leq \tilde{\mathbb{P}}[L>x] = 1 - G(x) \quad (x \in \mathbb{R}_+).\]
This corresponds to \textbf{first order stochastic dominance}.
The net risk premium with respect to \(G\) or \(\tilde{\mathbb{P}}\), respectively, satisfies
\begin{align*}
  \mathbb{E}_{\tilde{\mathbb{P}}}[L] &= \int_0^{\infty} (1-G(x))dx \\
    &= \int_0^{\infty} (1-F(x))dx + \int_0^{\infty} (F(x)-G(x))dx = \mathbb{E}_{\mathbb{P}}[L] + a
\end{align*}
with a safety loading \(a := \int_0^{\infty} (F(x)-G(x))dx > 0\) with respect to the actual net risk premium \(\mathbb{E}_{\mathbb{P}}[L]\).
In this sense, the actuary applies a change of measure before calculating the premium based on the expectation.


Let \(L \geq 0\) be a random variable which models possible claims of an insurance contract or a portfolio of insurance contracts (losses from the insurer's perspective).
Since monetary risk measures correspond to capital requirements, it is natural to define an insurance premium as follows:
\(\pi(L) = \rho(-L)\) for a single contract
\(\pi_n(L_i) = \frac{1}{n} \rho(-(L_1 + \dots + L_n))\) premium per risk in a homogenous portfolio of independent \(L_1, \dots, L_n\)
This change of sign is due to the actuarial interpretation: Here \(L\) describes the possible claims of an 'insurance risk'. 
The corresponding financial position of the insurer is then given by \(-L\), and so the premium \(\pi(L) = \rho(-L)\) is determined 
as the capital that is required to make that position acceptable.

Let \(\rho\) be a convex risk measure on \(\mathcal{X} = L^{\infty}(\Omega, \mathcal{F})\). 
We assume that \(\rho\) is finite and normalized to \(\rho(0)=0\).
In this case, the functional \(\pi\) on \(\mathcal{X}\) defined by
\[\pi(L) := \rho(-L)\]
satisfies the following properties:

Normalization: \(\pi(0)=0\),
Translation property: \(\pi(L+m) = \pi(L) + m\) for \(L \in \mathcal{X}\) and \(m \in \mathbb{R}\),
Monotonicity: \(L_1 \leq L_2 \Rightarrow \pi(L_1) \leq \pi(L_2)\),
Convexity: \(\pi\) is convex on \(\mathcal{X}\).

Note also that \(\pi\) is continuous from below if \(\rho\) is continuous from above.
Such functionals \(\pi\) were introduced by Deprez and Gerber (1985) in an actuarial context as convex principles of premium calculation.

The premium principle
\[\pi_{\gamma}(L) := \frac{1}{\gamma} \log E_{\mathbb{P}}[e^{\gamma L}],\]
which corresponds to the convex entropic risk measure \(e_{\gamma}\) is called the exponential principle.

The exponential principle \(\pi_{\gamma}\) satisfies the zero utility principle with respect to the exponential utility function 
\(u_{\gamma}(x) := \frac{1}{\gamma}(1-e^{-\gamma x})\), i.e.,
\[E_{\mathbb{P}}[u_{\gamma}(\pi_{\gamma}(L)-L)] = u_{\gamma}(0),\]
and it is additive on independent risks.

Conversely, any premium principle which satisfies the zero utility principle and is additive for independent risks must be of the form 
\(\pi_{\gamma}\) for some \(\gamma \geq 0\), including the linear case \(\gamma = 0\).

If \(L \geq 0\) models the possible claims of an insurance contract, the net risk premium is given by
\(\mathbb{E}_{\mathbb{P}}[L] = \int_0^{\infty} \mathbb{P}[L>x\) dx.
In order to obtain a safety loading, one could weight the probabilities \(\mathbb{P}[L>x]\) with a concave distortion function \(\psi\), i.e., an increasing and concave function on the unit interval defined by \(\psi(0)=0\) and \(\psi(1)=1\).
In this case, we obtain the premium
\[\pi(L) = \int_0^{\infty} \psi(\mathbb{P}[L>x]) dx\]
Thus, the premium is calculated as the Choquet integral
\[\pi(L) = \int L dc\]
with respect to some concave distortion \(c = \psi \circ \mathbb{P}\) of the underlying probability measure \(\mathbb{P}\).
In the actuarial literature this is often called Wang's premium principle. 
It corresponds to the class of distortion risk measures, up to a change of sign.
